% Nouveau EGA V (cote 25): archive cote 25, site page 18 (PDF page 19).
% typescript. Draft compared: none.
% AI-transcribed by Claude Opus 5.5 (claude-opus-5-5, adaptive thinking, effort high), Message Batches API, batch ega25b (msgbatch_01A7n5YuXubUPNK1zsepi8RG), 2026-11-13; not yet checked by eye.
\scanpage{25/18}
\mspage{15}

\inserted[typed]{sur $k(\hand{\eta})$.}
qui sont géom connexes en vertu de 3.1. , et \xout{\illegible{1}} l'intersection
de deux consécutifs étant $\neq \emptyset$\xout{\unsure{0}} en vertu de 2.3. Il s'ensuit (comme
$Y_{\hand{\eta}} = X_{\hand{\eta}}$ est la réunion des $X_{i\hand{\eta}}$ , $X_i$ parcourant l'ensemble des comp
irr de X ) que \xout{\illegible{1}} $Y_{\hand{\eta}}$ est géom connexe sur $k(\hand{\eta})$, cqfd.

2\unsure{0|D}.4. \uline{Etude d'une section hyperplane variable}: \hand{sections ``\uline{assez générales}''} .
4.1.
Nous revenons à la situation générale du nº 1, S étant un préschéma
quelconque. De plus, nous supposons X de présentation finie sur S.
De façon générale, notons que si P(Z,k) est une propriété''constructib-%
le'' \xout{\unsure{de préschéma sur}} d'un préschéma algébrique Z sur le corps k,
alors l'ensemble des $\hand{\zeta} \hand{\in} \overset{\hand{\vee}}{P}$ tels qu'on ait $P(Y_{\hand{\zeta}},k(\hand{\zeta}))$ est \uline{constructi-%
ble}, comme on voit en notant que $Y_{\hand{\zeta}}$ est la fibre en $\hand{\zeta}$ de $\underline{Y} \longrightarrow \overset{\hand{\vee}}{P}$\hand{,}qui
est un morphisme de prés finie, et en appliquant le par.9 . On a une
remarque analogue pour une\hand{|}propriété P(Z,M,k), où Z,k sont comme dessus
et M est un Module cohérent sur Z: si F est \deleted{cohérent}\inserted{de prés. finie} sur X, l'ensemble
des $\hand{\zeta} \hand{\in} \overset{\hand{\vee}}{P}$ tels qu'on ait $P(Y_{\hand{\zeta}},G_{\hand{\zeta}},k(\hand{\zeta}))$ est constructible. D'autre
part, dans les deux numéros suivants, on a développé \xout{\illegible{1}} dans divers
cas un critère pour que l'ensemble précédent E contienne le point
générique de $\overset{\hand{\vee}}{P}$ , S étant le spectre d'un corps k ; étantconstructible,
il en résulte que E contient un ouvert non vide: c'est le passage d'une
conclusion pour la \xout{x}section hyperplane générique à la conclusion analo\xout{g}
gue pour les sections hyperp\xout{h}lanes ``assez générales''. \xout{Enfin,} \hand{[}\deleted{Nous donne-%
rons également dans ce nº des critères permettant d'affirmer que E est
ouvert, en appliquant les résultats du par.12}\hand{]}. Noter d'ailleurs que \margin{$\downarrow$}
dans le cas S=Spec(k), on a également une réciproque: pour que la \xout{secti}secti-%
on hyperplane \uline{générique} satisfasse la propriété P, il f. et s. que
les $Y_{\hand{\zeta}}$ pour $\hand{\zeta}$ dans un voisinage convenable de $\hand{\eta}$ y satisfassent, et il
suffit de l'exiger pour $\hand{\zeta}$ \uline{fermé} dans $\overset{\hand{\vee}}{P}$ (ce qui, pour k algébriquement
clos, revient à envisager des points rationnels sur k, i.e. des

\margin{18}
